\section{讨论与结论}

本文的主结论是：在 API 成本约束下，代码生成系统的收益来自“把预算投向不同失败形态”，而不是统一增加调用次数。双基准主结果表明，同一协议下从单步到级联形成稳定前沿；中档主要提升来自分流与失败后修补，高档主要提升来自候选扩展与筛选组织。

机制证据与主结论需要分层阅读。LR+CV 在静态特征上接近随机，只能说明\textbf{离线失败标签}难以被表特征稳定预测；表~\ref{tab:route_random_multiset} 给出\textbf{固定 multiset 的整卷随机置换}对照，用于回答「在 hard 覆盖率不变时，结构—槽位对齐是否仍有边际」：MBPP 上 structural 较 random 通过率更高但 token 亦更高，HumanEval 上通过率相同而 token 差较小。该对照\textbf{不}等价于所有可能随机路由；结论应限定在本文实现与单次 seed 观测内。反馈形态多种子对照（表~\ref{tab:e2_feedback_format_5seed}）显示两组均值通过率差异落在各自标准差量级内，未做 formal 检验；在此证据强度下，更稳妥的部署表述是“compressed 回注在均值上未显示精度损失，同时平均 token 更低”，而非“默认工程形态必然更优”。

对部署实践而言，可形成清晰建议：中预算优先保留 \texttt{hard\_spec} 与失败回注修补，高预算将新增计算优先分配给 \texttt{extra} 候选扩展，再按需开启 refine 微调。这一策略与第~\ref{sec:cost_rho} 节 $\rho$ 前沿一致，即在可接受 token 增量下换取更高通过率。

当前证据边界也明确：细粒度消融仍以 MBPP 为主；尽管主表中的 MBPP 与 HumanEval 均已补齐四档三 seed 区间，multiset 随机路由对照仍为\textbf{单次 seed=42}，若需与主表同级的区间统计可扩展多 seed 重复配对。跨 API 版本漂移仍待系统覆盖。后续工作将优先扩展跨 API 版本稳定性验证，并在同一函数级代码生成口径下做更长周期回归。